%!TEX root = ../QuantumNoiseAndDecoherence.tex
% ──────────────────────────────────────────────────────────────
%  Lecture 10 — Fermionic Master Equation and Repeated Measurements
% ──────────────────────────────────────────────────────────────

% Continues Section 5 (Lindblad equations); opens Section 7
% (Classical noise and stochastic calculus).

\subsection{Quantum-dot master equation}\label{sec:dot-master-eq}

We return to the fermionic system introduced in the previous lecture:
a single quantum dot coupled to two external leads (source~$L$ and
drain~$R$) by the Hamiltonian
\begin{equation}\label{eq:dot-hamiltonian}
  H = \underbrace{%
    \sum_{k}\epsilon_k\, a_k^\adj a_k
    + E_c\, c^\adj c
    + \sum_{p}\epsilon_p\, d_p^\adj d_p}_{H_0}
  + \underbrace{%
    \sum_{k} T_k\, a_k^\adj c
    + \sum_{p} T_p^{\prime}\, c^\adj d_p
    + \text{h.c.}}_{V}\,,
\end{equation}
where $a_k$ annihilates a fermion in mode~$k$ of the source, $c$ is
the annihilation operator for the single bound state of the dot
(energy~$E_c$), and $d_p$ annihilates a fermion in mode~$p$ of the
drain.  An external voltage biases the leads so that electrons flow
from source to drain through the dot.

\begin{intuition}[Fermionic analogue of cavity decay]
  This model is the fermionic counterpart of the optical cavity
  coupled to input and output modes: electrons transit through the
  dot just as photons transit through a cavity.
\end{intuition}

\subsubsection{Interaction picture and reservoir correlators}

Going to the interaction picture with respect to~$H_0$ and expanding
the von~Neumann equation to second order---exactly as in the
Born--Markov derivation for cavity decay---the dynamics depend on two
types of non-vanishing reservoir correlators:
\begin{equation}\label{eq:correlators}
  C_n^{(L)}(\tau)
  = \sum_k \lvert T_k\rvert^2\,
    \bar{n}_k^{(L)}\, e^{-i(\epsilon_k - E_c)\tau}\,,
  \qquad
  C_a^{(L)}(\tau)
  = \sum_k \lvert T_k\rvert^2\,
    (1-\bar{n}_k^{(L)})\, e^{i(\epsilon_k - E_c)\tau}\,,
\end{equation}
with $\bar{n}_k^{(L)} = \expect{a_k^\adj a_k}_\beta$ the thermal
occupation of mode~$k$ in the source (analogous expressions hold for
the drain).  The anti-normally ordered correlator~$C_a$ is non-zero
because of fermionic anti-commutation relations---absent in the
bosonic case.

\subsubsection{Continuum limit and Markov approximation}

Replacing the discrete mode sum by an integral with density of
states~$\rho(\omega)$, only frequencies near~$E_c$ contribute
appreciably (stationary-phase argument).

\begin{keyresult}[Fermi--Dirac occupation]
  The mean occupation of a fermionic mode at energy~$\omega$ is
  \begin{equation}\label{eq:fermi-dirac}
    \eqbox{\bar{n}(\omega)
    = \frac{1}{e^{\beta(\omega - \omega_F)} + 1}}\,,
  \end{equation}
  where $\omega_F$ is the Fermi energy.  At low temperature,
  $\bar{n}(\omega) \to \Theta(\omega_F - \omega)$: all modes below
  the Fermi surface are filled.
\end{keyresult}

The bias sets $E_c < \omega_F^{(L)}$ and $E_c > \omega_F^{(R)}$;
at low temperature $\bar{n}^{(L)}(E_c) \approx 1$ and
$\bar{n}^{(R)}(E_c) \approx 0$.  The Markov approximation then
yields effective coupling rates
\begin{equation}\label{eq:gamma-rates}
  \frac{\gamma_L}{2}
  = \pi\,\rho(E_c)\,\lvert T(E_c)\rvert^2\,
    \bar{n}^{(L)}(E_c)\,,
  \qquad
  \frac{\gamma_R}{2}
  = \pi\,\rho(E_c)\,\lvert T'(E_c)\rvert^2\,
    \bigl(1 - \bar{n}^{(R)}(E_c)\bigr)\,.
\end{equation}

\subsubsection{Master equation for the dot}

Following the same steps as for cavity decay, the master equation for
the reduced state~$\rho$ of the dot ($\hilb \cong \C^2$) is
\begin{equation}\label{eq:dot-lindblad}
  \eqbox{\frac{\dd\rho}{\dd t}
  = \gamma_L\Bigl(c^\adj\rho\, c
      - \tfrac{1}{2}\{c\, c^\adj,\,\rho\}\Bigr)
  + \gamma_R\Bigl(c\,\rho\, c^\adj
      - \tfrac{1}{2}\{c^\adj c,\,\rho\}\Bigr)}\,.
\end{equation}
The first term describes electrons tunnelling \emph{onto} the dot
from the source (rate~$\gamma_L$); the second describes electrons
tunnelling \emph{off} the dot into the drain (rate~$\gamma_R$).

\begin{remark}
  Since the dot has one bound state and the electron is spin-polarised,
  the fermionic super-selection rule forbids superpositions of
  $\ket{0}$ and $\ket{1}$; the dynamics reduce to a single number.
\end{remark}

\subsubsection{Solution: occupation dynamics}

Defining $\expect{\hat{n}} = \tr(c^\adj c\,\rho)$ and using
fermionic identities ($c^2 = 0$, $(c^\adj c)^2 = c^\adj c$), one
obtains
\begin{equation}\label{eq:occupation-ode}
  \eqbox{\frac{\dd}{\dd t}\expect{\hat{n}}
  = \gamma_L\bigl(1 - \expect{\hat{n}}\bigr)
    - \gamma_R\,\expect{\hat{n}}}\,.
\end{equation}
The steady-state occupation is
$\expect{\hat{n}}_\infty = \gamma_L / (\gamma_L + \gamma_R)$,
and the full solution with initial condition
$\expect{\hat{n}}(0) = n_0$ reads
\begin{equation}\label{eq:n-solution}
  \expect{\hat{n}}(t)
  = \expect{\hat{n}}_\infty
    + \bigl(n_0 - \expect{\hat{n}}_\infty\bigr)\,
      e^{-(\gamma_L + \gamma_R)\,t}\,.
\end{equation}
Regardless of its initial state, the dot relaxes exponentially to
$\expect{\hat{n}}_\infty$ at rate $\gamma_L + \gamma_R$.

\begin{exercise}\label{ex:derive-dot-lindblad}
  Derive~\eqref{eq:dot-lindblad} by carrying out the Born--Markov
  procedure for~\eqref{eq:dot-hamiltonian}, replacing the bosonic bath
  correlators with~$\gamma_L/2$ and~$\gamma_R/2$.
\end{exercise}

% ──────────────────────────────────────────────────────────────
\subsection{Lindblad structure: a general pattern}%
\label{sec:lindblad-pattern}

The master equations derived so far---two-level atom, cavity decay,
quantum dot---all share the \emph{Lindblad form}:
\begin{equation}\label{eq:generic-lindblad}
  \frac{\dd\rho}{\dd t}
  = -i[H_{\mathrm{sys}},\,\rho]
    + \sum_j \gamma_j\,\mathcal{D}[L_j]\,\rho\,,
  \qquad
  \mathcal{D}[L]\,\rho
  = L\rho L^\adj
    - \tfrac{1}{2}\{L^\adj L,\,\rho\}\,.
\end{equation}

\begin{intuition}[Why always Lindblad form?]
  Complete positivity of the evolved state for all~$t \geq 0$, under
  the Markov assumption, severely constrains the generator.  The
  Lindblad theorem (forthcoming) shows
  that~\eqref{eq:generic-lindblad} is the \emph{most general} such
  generator.
\end{intuition}

% ──────────────────────────────────────────────────────────────
\section{Classical noise and stochastic calculus}%
\label{sec:classical-noise}

\subsection{Repeated quantum measurements}%
\label{sec:repeated-measurements}

We now present a derivation that provides perhaps the clearest
physical intuition for the Lindblad structure.  It arises from a model
of \emph{repeated quantum measurements} due to Caves and
Milburn~\cite{CavesMilburn87}.

\begin{historical}[Caves--Milburn (1987)]
  In the Caves--Milburn model the Lindblad master equation arises
  \emph{exactly}---not as an approximation---from repeatedly coupling
  a system to independent meters.
\end{historical}

\subsubsection{The setup}

Consider a system~$S$ with Hamiltonian~$H_{\mathrm{sys}}$ and
$n$~independent meters $M_1, \ldots, M_n$ (free particles).  The
total Hilbert space is
\begin{equation}\label{eq:total-hilb}
  \hilb_{\mathrm{tot}}
  = \hilb_S \tp \hilb_{M_1} \tp \cdots \tp \hilb_{M_n}\,.
\end{equation}
The system evolves freely except at discrete times $t_j = j\tau$,
when it interacts instantaneously with meter~$M_j$ via the
\textbf{von~Neumann measurement Hamiltonian}
\begin{equation}\label{eq:vn-hamiltonian}
  \eqbox{H_{\mathrm{int}}^{(j)}
  = \kappa\, \hat{A} \tp \hat{p}_j}\,,
\end{equation}
where $\hat{A}$ is the system observable to be measured, $\hat{p}_j$
is the momentum of meter~$j$, and $\kappa$ controls the interaction
strength.

\subsubsection{Von Neumann measurement prescription}

\begin{definition}[Von Neumann measurement model]%
\label{def:vn-measurement}
  The meter is prepared in the position eigenstate
  $\ket{x = 0}_{M}$.  The system--meter composite evolves under
  $H_{\mathrm{int}}$ for a time~$\tau$, producing the unitary
  $U = \exp(-i\kappa\tau\, \hat{A} \tp \hat{p})$.
\end{definition}

Expanding $\ket{\psi} = \sum_j c_j\ket{\alpha_j}$ in the eigenbasis
of~$\hat{A}$, the initial product state
$\ket{\Psi(0)} = \ket{\psi} \tp \ket{x=0}_M$ evolves to
\begin{equation}\label{eq:posterior-entangled}
  \eqbox{\ket{\Psi(\tau)}
  = \sum_j c_j\,\ket{\alpha_j}
    \tp \ket{x = \kappa\tau\alpha_j}_M}\,,
\end{equation}
since $e^{-i\kappa\tau\alpha_j\hat{p}}$ is the translation operator
shifting position by~$\kappa\tau\alpha_j$.  Each eigenvalue is now
correlated with a distinct pointer reading.

\begin{keyresult}[Effective projective measurement]
  A subsequent projective measurement of the pointer position yields
  outcome $x = \kappa\tau\alpha_j$ with probability~$|c_j|^2$, and
  the system collapses to~$\ket{\alpha_j}$.  This reproduces the
  textbook measurement postulate for~$\hat{A}$.
\end{keyresult}

\subsubsection{Repeated measurements and scaling regimes}

The full time-dependent Hamiltonian is
\begin{equation}\label{eq:repeated-H}
  H(t) = H_{\mathrm{sys}}
  + \kappa \sum_{j=1}^{n}
    \delta(t - j\tau)\;\hat{A} \tp \hat{p}_j\,.
\end{equation}
The physics depends on the scaling of~$n$ and~$\kappa$:
\begin{itemize}
  \item \textbf{Strong coupling} ($\kappa$ large): each interaction
    is nearly projective.  Increasing the measurement rate freezes the
    system in an eigenstate of~$\hat{A}$---the \textbf{quantum Zeno
    effect}.
  \item \textbf{Weak coupling} ($\kappa \to 0$, $n$ fixed): the
    meters decouple; the system evolves unitarily under
    $H_{\mathrm{sys}}$.
  \item \textbf{Intermediate scaling} ($n \to \infty$,
    $\kappa \to 0$, $n\kappa^2$ fixed): the meters exert a finite
    back-action.  This yields a Lindblad master equation and provides
    a transparent model for environmental decoherence.
\end{itemize}

\begin{intuition}[The Goldilocks regime]
  Each individual measurement barely disturbs the system, but
  infinitely many weak measurements produce continuous, irreversible
  evolution---precisely how a thermal environment induces decoherence.
\end{intuition}

The derivation of the Lindblad equation from this
repeated-measurement model will be carried out in the next lecture.
